\section{SSH részletesen: kulcsok, konfiguráció és jogosultságok}

Az SSH aszimmetrikus kulcsokkal biztosít titkosított kommunikációt. A publikus kulcs a szerveren kerül elhelyezésre, a privát kulcs a kliensnél marad.

\subsection{Kulcsgenerálás és titkosítás}
\begin{lstlisting}[language=bash]
# RSA 4096 bites kulcs
ssh-keygen -t rsa -b 4096 -f ~/.ssh/id_rsa
# Ed25519 (smaller and faster)
ssh-keygen -t ed25519 -a 100 -f ~/.ssh/id_ed25519
# Protect the private key with a passphrase
ssh-keygen -p -f ~/.ssh/id_rsa
\end{lstlisting}

\subsection{authorized\_keys és jogosultságok}
\begin{itemize}
  \item A \path{~/.ssh/authorized_keys} fájlban tárolt publikuskulcsok határozzák meg a bejelentkezést.
  \item Jogosultságok: \path{chmod 700 ~/.ssh} és \path{chmod 600 ~/.ssh/authorized_keys}.
\end{itemize}

\subsection{SSH másolása és automatizálás}
\begin{lstlisting}[language=bash]
# Copy the public key to the server
ssh-copy-id user@host
# Use a specific key
ssh -i ~/.ssh/id_ed25519 user@host
\end{lstlisting}

Biztonsági megjegyzés:
\begin{itemize}
  \item Ne osszuk meg a privát kulcsot.
  \item GPG-vel vagy hardware tokennel (YubiKey) tárolt kulcs további védelmet ad.
\end{itemize}
